% Worked practice problems for midterm 1, set in the style of my lecture scribes.
% The style files are the ones in ../scribes (scribe.sty is the unmodified course file).
% This document is not part of the scribes. Build from this folder with
%   pdflatex -interaction=nonstopmode practiceproblems.tex   (three passes)

\documentclass[a4paper,11pt,final]{article}
\makeatletter
\def\input@path{{../scribes/}}
\makeatother
\usepackage{scribecompat} % must come before scribe
\usepackage{scribe}
\usepackage{scribeextras} % dashed exercise boxes, Dirac notation, inline TikZ

\theoremstyle{plain}
\newtheorem{claim}[theorem]{Claim}
\crefname{claim}{Claim}{Claims}
\Crefname{claim}{Claim}{Claims}

% Practice problems are numbered inside their section: Problem 1.1, 1.2, ...
\theoremstyle{definition}
\newtheorem{pproblem}{Problem}[section]
\crefname{pproblem}{Problem}{Problems}
\Crefname{pproblem}{Problem}{Problems}

% The table of contents lists the sections with their numbers and page numbers
% (scribeextras.sty hides both for the scribes, where the lecture days carry them).
\makeatletter
\renewcommand{\cftsecpresnum}{}
\renewcommand{\cftsecaftersnum}{}
\setlength{\cftsecnumwidth}{2.3em}
\setlength{\cftsecindent}{0em}
\renewcommand{\cftsecfillnum}[1]{{\cftsecleader}\nobreak
    \makebox[\@pnumwidth][\cftpnumalign]{\cftsecpagefont #1}\cftsecafterpnum\par}
\makeatother

% A little space between the parts (a), (b), ... of a solution.
\AtBeginEnvironment{wsolution}{\setlength{\parskip}{0.35em}}

\begin{document}

\begin{center}
\textbf{\large Quantum Computing} \\

Professor Gopal Pandurangan\\
\vspace{0.1in}
\end{center}

\hrule

\vspace{0.1in}

\textbf{Practice Problems for Midterm 1} (lectures of August 25 -- October 8, 2026)

\vspace{0.1in}

\hrule

\vspace{0.1in}

These are practice problems in the style of the likely question types, each with a full solution.
They are a guess at the exam, not questions from the professor. Every method used is from the
scribes; the lecture date is given to find the matching page.
\Cref{sec:chernoff,sec:maxload} fill two gaps in the scribes: the proof of the Chernoff bound, and
the upper-tail version of the max-load exercise.

\tableofcontents

\clearpage

\section{Run Deutsch--Jozsa, Bernstein--Vazirani or Simon on a Concrete $f$}\label{sec:run}

The method is always the same: write the state after every gate. Two facts from the scribes
(Sep 24, Sep 29) do all the work:
\[U_f \ket{x}\ket{-} = (-1)^{f(x)} \ket{x}\ket{-},
    \qquad
    H^{\otimes n}\ket{x} = \frac{1}{\sqrt{2^n}} \sum_{y \in \set{0,1}^n} (-1)^{x \cdot y} \ket{y}.\]

\begin{exercisebox}
\begin{pproblem}[Deutsch--Jozsa with $n = 2$]\label{pp:dj}
    Let $f(x_1 x_2) = x_1 \oplus x_2$. Run the Deutsch--Jozsa algorithm on $f$. What is measured,
    and what does the algorithm conclude?
\end{pproblem}
\tcblower
\begin{wsolution}
    \leavevmode
    \begin{enumerate}
        \item Start in $\ket{00}\ket{1}$. Hadamards on all three qubits give
              $\frac{1}{2}\left(\ket{00} + \ket{01} + \ket{10} + \ket{11}\right)\ket{-}$.
        \item The query multiplies each $\ket{x}$ by $(-1)^{f(x)}$. Here $f(00) = 0$, $f(01) = 1$,
              $f(10) = 1$, $f(11) = 0$, so the state is
              $\frac{1}{2}\left(\ket{00} - \ket{01} - \ket{10} + \ket{11}\right)\ket{-}$.
        \item Hadamards on the first two qubits. The amplitude of $\ket{y}$ is
              $\frac{1}{4} \sum_{x} (-1)^{f(x) + x \cdot y}$.
        \item For $y = 00$ this is $\frac{1}{4}(1 - 1 - 1 + 1) = 0$. For $y = 11$ we have
              $x \cdot y = x_1 \oplus x_2 = f(x)$, so every term is $+1$ and the amplitude is $1$.
        \item The measurement gives $11$ with probability $1$. The outcome is not $00$, so $f$ is
              balanced.
    \end{enumerate}
    One quantum query was used. A deterministic classical algorithm needs $2^{n-1} + 1 = 3$ queries.

    \emph{Check with a constant function.} For $f \equiv 1$ the state after the query is
    $-\frac{1}{2} \sum_{x} \ket{x} = -\ket{+}\ket{+}$, and the Hadamards turn it into $-\ket{00}$.
    The measurement gives $00$, so $f$ is constant. The global sign $-1$ does not affect the
    measurement.
\end{wsolution}
\end{exercisebox}

\begin{exercisebox}
\begin{pproblem}[Bernstein--Vazirani with $n = 3$]\label{pp:bv}
    Let $f(x) = s \cdot x \bmod 2$ with the hidden string $s = 101$, so $f(x) = x_1 \oplus x_3$.
    \begin{enumerate}[label=(\alph*)]
        \item How does a classical algorithm find $s$?
        \item Run the quantum algorithm.
    \end{enumerate}
\end{pproblem}
\tcblower
\begin{wsolution}
    (a) Query the unit vectors: $f(100) = 1$, $f(010) = 0$, $f(001) = 1$. The answers are the bits
    of $s$, so $n = 3$ queries. Fewer cannot work, since each answer is one bit and $s$ has $n$ bits.

    (b) After the first Hadamards and the query, the first register is
    \[\frac{1}{\sqrt{8}} \sum_{x} (-1)^{s \cdot x} \ket{x}
        = \frac{1}{\sqrt{8}} \left(\ket{000} - \ket{001} + \ket{010} - \ket{011} - \ket{100} + \ket{101} - \ket{110} + \ket{111}\right).\]
    After the second Hadamards the amplitude of $\ket{y}$ is
    \[\frac{1}{8} \sum_{x} (-1)^{s \cdot x + y \cdot x} = \frac{1}{8} \sum_{x} (-1)^{(s \oplus y) \cdot x}
        = \begin{cases} 1 & \text{if } y = s, \\ 0 & \text{if } y \neq s, \end{cases}\]
    by the sum-of-signs claim (Sep 29): the sum is $2^n$ if $s \oplus y = 000$ and $0$ otherwise.
    So the measurement gives $y = 101 = s$ with probability $1$, after one query.
\end{wsolution}
\end{exercisebox}

\begin{exercisebox}
\begin{pproblem}[Simon with $n = 3$]\label{pp:simon}
    The function $f \colon \set{0,1}^3 \to \set{0,1}^3$ is given by the table below. Run Simon's
    algorithm and find $s$.
    \begin{center}
        \begin{tabular}{c|cccccccc}
            $x$    & 000 & 001 & 010 & 011 & 100 & 101 & 110 & 111 \\ \hline
            $f(x)$ & 101 & 010 & 000 & 110 & 000 & 110 & 101 & 010
        \end{tabular}
    \end{center}
\end{pproblem}
\tcblower
\begin{wsolution}
    \leavevmode
    \begin{enumerate}
        \item After the Hadamards and the query, the state is
              $\frac{1}{\sqrt{8}} \sum_{x} \ket{x}\ket{f(x)}$.
        \item Measure the second register. Suppose the outcome is $010$. The inputs with
              $f(x) = 010$ are $001$ and $111$, so the first register is
              $\frac{1}{\sqrt{2}}\left(\ket{001} + \ket{111}\right)$.
        \item Apply $H^{\otimes 3}$. The amplitude of $\ket{y}$ is
              \[\frac{1}{\sqrt{2}} \cdot \frac{1}{\sqrt{8}} \left((-1)^{001 \cdot y} + (-1)^{111 \cdot y}\right)
                  = \frac{1}{4} (-1)^{001 \cdot y} \left(1 + (-1)^{110 \cdot y}\right).\]
        \item This is $\pm\frac{1}{2}$ if $110 \cdot y = 0$ and $0$ if $110 \cdot y = 1$. So $y$ is
              one of $000, 001, 110, 111$, each with probability $\sfrac{1}{4}$. The state is
              $\frac{1}{2}\left(\ket{000} - \ket{001} + \ket{110} - \ket{111}\right)$.
        \item Every outcome satisfies $y_1 \oplus y_2 = 0$. The same holds for any other outcome in
              step 2, because every pair $\set{x, x \oplus s}$ has the same difference.
        \item Run the circuit $n - 1 = 2$ times. Suppose the outcomes are $y = 111$ and $y = 001$.
              They give the equations $s_1 \oplus s_2 \oplus s_3 = 0$ and $s_3 = 0$.
        \item Solve: $s_3 = 0$ and $s_1 = s_2$. The solutions are $000$ and $110$, so the nonzero
              candidate is $s' = 110$.
        \item Check with two classical queries: $f(000) = 101 = f(110)$. So $f$ is $2$-to-$1$ and
              $s = 110$.
    \end{enumerate}

    \emph{What can go wrong.} If the two outcomes are dependent (for example $000$ and $111$, or
    $001$ twice), the system has more than two solutions and we run again. Two outcomes are
    independent with probability $\frac{3}{4} \cdot \frac{1}{2} = \frac{3}{8}$: the first must be
    nonzero ($3$ of $4$), and the second must be outside $\set{000, y_1}$ ($2$ of $4$). This matches
    the bound of at least $\sfrac{1}{4}$ from Oct 1.

    \emph{If $f$ were $1$-to-$1$.} Step 2 would leave a single $\ket{x}$, every $y$ would have
    probability $\sfrac{1}{8}$, and the check in step 8 would fail, so the algorithm would answer
    $s = 000$.
\end{wsolution}
\end{exercisebox}

\section{Grover with Given $N$ and $a$}\label{sec:grover}

Everything follows from three formulas (Oct 6, Oct 8). $N$ is the number of inputs, $a$ the number
of solutions, and $k$ the number of iterations:
\[\sin\theta = \sqrt{\frac{a}{N}},
    \qquad
    \prob{\text{success after } k \text{ iterations}} = \sin^2\left((2k+1)\theta\right),
    \qquad
    k = \floor{\frac{\pi}{4\theta}}.\]
The recipe: compute $\theta$, compute $k$, plug in. Each iteration adds $2\theta$ to the angle, and
we want the angle $(2k+1)\theta$ as close to $90^\circ$ as possible.

\begin{exercisebox}
\begin{pproblem}[$N = 4$, $a = 1$]\label{pp:grover4}
    Find $\theta$, the best $k$ and the success probability. Then verify by tracking the four
    amplitudes.
\end{pproblem}
\tcblower
\begin{wsolution}
    \leavevmode
    \begin{enumerate}
        \item $\sin\theta = \sqrt{\sfrac{1}{4}} = \frac{1}{2}$, so $\theta = 30^\circ = \sfrac{\pi}{6}$.
        \item $k = \floor{\sfrac{\pi}{(4\theta)}} = \floor{1.5} = 1$.
        \item The success probability is $\sin^2(3 \cdot 30^\circ) = \sin^2 90^\circ = 1$. One query
              finds the solution with certainty; classically $3$ queries are needed in the worst
              case.
    \end{enumerate}

    \emph{Amplitude check.} Start with $\left(\frac{1}{2}, \frac{1}{2}, \frac{1}{2}, \frac{1}{2}\right)$.
    The phase flip $Z_f$ negates the solution:
    $\left(\frac{1}{2}, \frac{1}{2}, \frac{1}{2}, -\frac{1}{2}\right)$, with the solution last. The
    mean is $\frac{1}{4}\left(\frac{1}{2} + \frac{1}{2} + \frac{1}{2} - \frac{1}{2}\right) = \frac{1}{4}$.
    Inversion about the mean sends each amplitude $\alpha$ to $2 \cdot \frac{1}{4} - \alpha$, which
    gives $(0, 0, 0, 1)$.
\end{wsolution}
\end{exercisebox}

\begin{exercisebox}
\begin{pproblem}[$N = 8$, $a = 1$, and overshooting]\label{pp:grover8}
    Find the best $k$, and the success probability for $k = 1, 2, 3$.
\end{pproblem}
\tcblower
\begin{wsolution}
    \leavevmode
    \begin{enumerate}
        \item $\sin\theta = \sfrac{1}{\sqrt{8}} \approx 0.3536$, so
              $\theta \approx 0.3614$ rad $\approx 20.7^\circ$.
        \item $k = \floor{\sfrac{\pi}{(4 \cdot 0.3614)}} = \floor{2.17} = 2$.
        \item Plug in:
    \end{enumerate}
    \begin{center}
        \begin{tabular}{ccc}
            \toprule
            $k$ & Angle $(2k+1)\theta$ & Success probability                \\ \midrule
            $0$ & $20.7^\circ$         & $\sfrac{1}{8} = 0.125$             \\
            $1$ & $62.1^\circ$         & $\sfrac{25}{32} \approx 0.781$     \\
            $2$ & $103.5^\circ$        & $\sfrac{121}{128} \approx 0.945$   \\
            $3$ & $144.9^\circ$        & $\sfrac{169}{512} \approx 0.330$   \\ \bottomrule
        \end{tabular}
    \end{center}
    The best is $k = 2$ with probability about $0.945$. A third iteration rotates past $90^\circ$
    and the probability drops to $0.33$: more iterations are not better.
\end{wsolution}
\end{exercisebox}

\begin{exercisebox}
\begin{pproblem}[$N = 8$, $a = 2$]\label{pp:grover82}
    Same $N$, two solutions: find $\theta$, $k$ and the success probability.
\end{pproblem}
\tcblower
\begin{wsolution}
    $\sin\theta = \sqrt{\sfrac{2}{8}} = \frac{1}{2}$, so $\theta = 30^\circ$, $k = 1$ and the
    success probability is $\sin^2 90^\circ = 1$. Only $\sfrac{a}{N}$ matters, so this is the same as
    $N = 4$, $a = 1$. Whenever $a = \sfrac{N}{4}$, one iteration succeeds with certainty.
\end{wsolution}
\end{exercisebox}

\begin{exercisebox}
\begin{pproblem}[Using the wrong $a$]\label{pp:wrong-a}
    Let $N = 64$. We run the number of iterations that is right for $a = 1$, but in fact $a = 4$.
    What is the success probability? What should $k$ have been?
\end{pproblem}
\tcblower
\begin{wsolution}
    \leavevmode
    \begin{enumerate}
        \item For $a = 1$: $\sin\theta = \sfrac{1}{8}$, $\theta \approx 0.1253$ and
              $k = \floor{6.27} = 6$. With $a = 1$ this would succeed with probability
              $\sin^2(13 \cdot 0.1253) \approx 0.997$.
        \item With $a = 4$: $\sin\theta = \sqrt{\sfrac{4}{64}} = \frac{1}{4}$ and
              $\theta \approx 0.2527$. After $6$ iterations the angle is
              $13 \cdot 0.2527 \approx 3.285$ rad $\approx 188^\circ$, and $\sin^2 \approx 0.02$.
        \item The right choice is $k = \floor{\sfrac{\pi}{(4 \cdot 0.2527)}} = \floor{3.11} = 3$, with
              success probability $\sin^2(7 \cdot 0.2527) \approx 0.961$.
    \end{enumerate}
    So with the wrong $a$ the algorithm fails $98\%$ of the time. This is why an unknown $a$ needs a
    random $k$ (Oct 8): choose $k$ uniformly from $\set{1, \dots, \floor{2\sqrt{N}}}$ and the success
    probability is at least $\sfrac{3}{8}$ for every $a$ with $1 \leq a \leq \sfrac{N}{2}$
    (\cref{pp:unknown-a}).
\end{wsolution}
\end{exercisebox}

\begin{exercisebox}
\begin{pproblem}[$a = N/2$ and $a > N/2$]\label{pp:half}
    What does Grover's algorithm do when half or more of the inputs are solutions?
\end{pproblem}
\tcblower
\begin{wsolution}
    For $a = \sfrac{N}{2}$: $\theta = 45^\circ$, and the angles $45^\circ, 135^\circ, 225^\circ, \dots$
    all have $\sin^2 = \frac{1}{2}$. Iterating never helps. For $a > \sfrac{N}{2}$:
    $\theta > 45^\circ$, so $k = \floor{\sfrac{\pi}{(4\theta)}} = 0$, and measuring the uniform state
    directly succeeds with probability $\sfrac{a}{N} > \frac{1}{2}$. In both cases just sample a
    random $x$ and check it with one query.
\end{wsolution}
\end{exercisebox}

\section{Amplitude Amplification for a Procedure with Success Probability $p$}\label{sec:ampamp}

Amplitude amplification is Grover with the Hadamards replaced by any measurement-free algorithm
$\mathcal{A}$ (Oct 8). If $\mathcal{A}$ succeeds with probability $p$, then
\begin{gather*}
    \sin\theta = \sqrt{p},
    \qquad
    Q = -\mathcal{A}\, Z_0\, \mathcal{A}^{-1} Z_f, \\
    \prob{\text{success after } k \text{ rounds}} = \sin^2\left((2k+1)\theta\right),
    \qquad
    k = \floor{\frac{\pi}{4\theta}}.
\end{gather*}
The cost is $2k + 1$ applications of $\mathcal{A}$ or its inverse: one to prepare the state and two
per round.

\begin{exercisebox}
\begin{pproblem}[$p = 1/100$]\label{pp:amp100}
    A quantum algorithm $\mathcal{A}$ outputs a good answer with probability $\sfrac{1}{100}$, and we
    can check an answer with $f$. How many runs of $\mathcal{A}$ are needed classically, and how
    many with amplitude amplification?
\end{pproblem}
\tcblower
\begin{wsolution}
    \leavevmode
    \begin{enumerate}
        \item Classically: run, check, repeat. The number of runs is geometric with mean
              $\sfrac{1}{p} = 100$.
        \item $\sin\theta = \sqrt{\sfrac{1}{100}} = 0.1$, so $\theta \approx 0.1002$.
        \item $k = \floor{\sfrac{\pi}{(4 \cdot 0.1002)}} = \floor{7.84} = 7$.
        \item The success probability is $\sin^2(15 \cdot 0.1002) = \sin^2(1.503) \approx 0.995$.
        \item Cost: $2 \cdot 7 + 1 = 15$ applications of $\mathcal{A}$ or its inverse, and $7$
              queries to $f$.
    \end{enumerate}
    So $15$ applications instead of $100$. In general the cost is at most
    $\frac{\pi}{2\sqrt{p}} + 1 = \bigO{\sfrac{1}{\sqrt{p}}}$ against $\sfrac{1}{p}$.
\end{wsolution}
\end{exercisebox}

\begin{exercisebox}
\begin{pproblem}[$p = 1/4$]\label{pp:amp4}
    Answer the same question for $p = \sfrac{1}{4}$.
\end{pproblem}
\tcblower
\begin{wsolution}
    $\sin\theta = \frac{1}{2}$, $\theta = 30^\circ$, $k = 1$, and the success probability is
    $\sin^2 90^\circ = 1$. Three applications of $\mathcal{A}$ or its inverse give a good answer
    with certainty. Classically the expected number of runs is $4$, and no fixed number of runs
    gives certainty.
\end{wsolution}
\end{exercisebox}

\begin{exercisebox}
\begin{pproblem}[Short questions]\label{pp:ampshort}
    \leavevmode
    \begin{enumerate}[label=(\alph*)]
        \item Why must $\mathcal{A}$ contain no measurement?
        \item How is Grover's algorithm a special case?
        \item What if $p$ is unknown?
    \end{enumerate}
\end{pproblem}
\tcblower
\begin{wsolution}
    (a) The operator $Q$ uses the inverse of $\mathcal{A}$. A unitary circuit is inverted by running
    it backwards with each gate inverted. A measurement is not unitary and cannot be undone.

    (b) Take $\mathcal{A} = H^{\otimes n}$. Then $\mathcal{A}\ket{0^n}$ is the uniform state,
    $p = \sfrac{a}{N}$, and $Q = -H^{\otimes n} Z_0 H^{\otimes n} Z_f$ is the Grover iteration.

    (c) The formula for $k$ needs $p$. The theorem in the scribe assumes that $p$ is known; the
    unknown case was not covered.
\end{wsolution}
\end{exercisebox}

\section{Entangled or Product; Unitary or Not}\label{sec:entangled}

\paragraph{The test for two qubits.}
Write the state as $a\ket{00} + b\ket{01} + c\ket{10} + d\ket{11}$. Suppose it is a product
$\left(\alpha_0\ket{0} + \alpha_1\ket{1}\right)\left(\beta_0\ket{0} + \beta_1\ket{1}\right)$.
Multiplying out gives $a = \alpha_0\beta_0$, $b = \alpha_0\beta_1$, $c = \alpha_1\beta_0$,
$d = \alpha_1\beta_1$, and therefore
\[ad = \alpha_0\alpha_1\beta_0\beta_1 = bc.\]
So if $ad \neq bc$ the state is entangled. This is the argument for the Bell state (Sep 15) written
once for all states. To show that a state is a product, write down the two factors.

\begin{exercisebox}
\begin{pproblem}[Five states]\label{pp:five}
    Which of these states are entangled? Prove each answer.
    \begin{enumerate}[label=(\alph*)]
        \item $\frac{1}{2}\left(\ket{00} + \ket{01} + \ket{10} + \ket{11}\right)$
        \item $\frac{1}{2}\left(\ket{00} - \ket{01} + \ket{10} - \ket{11}\right)$
        \item $\frac{1}{2}\left(\ket{00} + \ket{01} + \ket{10} - \ket{11}\right)$
        \item $\frac{1}{\sqrt{2}}\left(\ket{01} - \ket{10}\right)$
        \item $\sqrt{\sfrac{1}{3}}\, \ket{00} + \sqrt{\sfrac{2}{3}}\, \ket{01}$
    \end{enumerate}
\end{pproblem}
\tcblower
\begin{wsolution}
    Compare $ad$ with $bc$:
    \begin{center}
        \begin{tabular}{cccl}
            \toprule
            State & $ad$              & $bc$              & Answer                                                                      \\ \midrule
            (a)   & $\sfrac{1}{4}$    & $\sfrac{1}{4}$    & Product: $\ket{+}\ket{+}$                                                   \\
            (b)   & $-\sfrac{1}{4}$   & $-\sfrac{1}{4}$   & Product: $\ket{+}\ket{-}$                                                   \\
            (c)   & $-\sfrac{1}{4}$   & $\sfrac{1}{4}$    & Entangled                                                                   \\
            (d)   & $0$               & $-\sfrac{1}{2}$   & Entangled                                                                   \\
            (e)   & $0$               & $0$               & Product: $\ket{0}\left(\sqrt{\sfrac{1}{3}}\,\ket{0} + \sqrt{\sfrac{2}{3}}\,\ket{1}\right)$ \\ \bottomrule
        \end{tabular}
    \end{center}

    \emph{Full write-up for (c).} Suppose the state equals
    $\left(\alpha_0\ket{0} + \alpha_1\ket{1}\right)\left(\beta_0\ket{0} + \beta_1\ket{1}\right)$.
    Then $\alpha_0\beta_0 = \frac{1}{2}$, $\alpha_0\beta_1 = \frac{1}{2}$,
    $\alpha_1\beta_0 = \frac{1}{2}$ and $\alpha_1\beta_1 = -\frac{1}{2}$. The product of the first
    and the last is $\alpha_0\alpha_1\beta_0\beta_1 = -\frac{1}{4}$. The product of the middle two
    is $\alpha_0\alpha_1\beta_0\beta_1 = +\frac{1}{4}$. This is a contradiction, so the state is
    entangled.
\end{wsolution}
\end{exercisebox}

\begin{exercisebox}
\begin{pproblem}[Is the matrix unitary?]\label{pp:unitary}
    Decide for each matrix whether it is unitary.
    \[U = \frac{1}{\sqrt{2}} \begin{pmatrix} 1 & i \\ i & 1 \end{pmatrix},
        \qquad
        V = \begin{pmatrix} 1 & 1 \\ 0 & 1 \end{pmatrix},
        \qquad
        W = \frac{1}{\sqrt{2}} \begin{pmatrix} 1 & 1 \\ 1 & 1 \end{pmatrix}.\]
\end{pproblem}
\tcblower
\begin{wsolution}
    Compute $U^\dagger U$, where $\dagger$ means: transpose, and conjugate every entry.
    \[U^\dagger U = \frac{1}{2} \begin{pmatrix} 1 & -i \\ -i & 1 \end{pmatrix} \begin{pmatrix} 1 & i \\ i & 1 \end{pmatrix}
        = \frac{1}{2} \begin{pmatrix} 1 - i^2 & i - i \\ -i + i & -i^2 + 1 \end{pmatrix}
        = \begin{pmatrix} 1 & 0 \\ 0 & 1 \end{pmatrix}.\]
    So $U$ is unitary. Do not forget the conjugate: without it the diagonal would be $1 + i^2 = 0$.

    $V$ is not unitary: $V\ket{1} = \ket{0} + \ket{1}$ has norm $\sqrt{2}$, and a unitary matrix
    preserves norms (Sep 15). $W$ is not unitary: $W\ket{0} = W\ket{1} = \ket{+}$, so two orthogonal
    states are sent to the same state; also $W\left(\ket{0} - \ket{1}\right) = 0$.

    \emph{Quick test.} A matrix is unitary exactly when its columns are orthonormal, with the
    complex inner product.
\end{wsolution}
\end{exercisebox}

\begin{exercisebox}
\begin{pproblem}[A small circuit]\label{pp:circuit}
    Apply $H$ to the first qubit and then CNOT (the first qubit is the control) to $\ket{10}$. Is
    the result entangled?
\end{pproblem}
\tcblower
\begin{wsolution}
    \leavevmode
    \begin{enumerate}
        \item $H$ on the first qubit: $\ket{10}$ becomes
              $\frac{1}{\sqrt{2}}\left(\ket{0} - \ket{1}\right) \otimes \ket{0}
                  = \frac{1}{\sqrt{2}}\left(\ket{00} - \ket{10}\right)$.
        \item CNOT flips the second qubit when the first is $1$:
              $\frac{1}{\sqrt{2}}\left(\ket{00} - \ket{11}\right)$.
        \item Test: $ad = -\frac{1}{2}$ and $bc = 0$, so the state is entangled.
    \end{enumerate}
    The same circuit sends $\ket{00}$ to $\frac{1}{\sqrt{2}}\left(\ket{00} + \ket{11}\right)$,
    $\ket{01}$ to $\frac{1}{\sqrt{2}}\left(\ket{01} + \ket{10}\right)$ and $\ket{11}$ to
    $\frac{1}{\sqrt{2}}\left(\ket{01} - \ket{10}\right)$.
\end{wsolution}
\end{exercisebox}

\section{Random Walk on a Small Graph}\label{sec:walk}

The recipe (Sep 8, Sep 10), for a graph with $m$ edges:
\[P_{uv} = \begin{cases} \dfrac{1}{\deg(u)} & \text{if } u \sim v, \\[1ex] 0 & \text{otherwise,} \end{cases}
    \qquad
    p_{t+1} = p_t P,
    \qquad
    \pi_v = \frac{\deg(v)}{2m}.\]
The walk converges to $\pi$ from every start exactly when the graph is connected and not bipartite.
A graph is bipartite exactly when it has no odd cycle.

\begin{exercisebox}
\begin{pproblem}[A triangle with a tail]\label{pp:triangle}
    Vertices $1, 2, 3$ form a triangle and vertex $4$ is joined to vertex $3$.
    \begin{enumerate}[label=(\alph*)]
        \item Write $P$.
        \item Find $\pi$ and verify it.
        \item Does the walk converge?
        \item Starting at vertex $4$, compute $p_1, p_2, p_3$.
    \end{enumerate}
\end{pproblem}
\tcblower
\begin{wsolution}
    (a) The degrees are $2, 2, 3, 1$, so
    \[P = \begin{pmatrix}
            0            & \tfrac{1}{2} & \tfrac{1}{2} & 0            \\
            \tfrac{1}{2} & 0            & \tfrac{1}{2} & 0            \\
            \tfrac{1}{3} & \tfrac{1}{3} & 0            & \tfrac{1}{3} \\
            0            & 0            & 1            & 0
        \end{pmatrix}.\]
    Every row sums to $1$.

    (b) There are $m = 4$ edges, so
    $\pi = \left(\frac{2}{8}, \frac{2}{8}, \frac{3}{8}, \frac{1}{8}\right)
        = \left(\frac{1}{4}, \frac{1}{4}, \frac{3}{8}, \frac{1}{8}\right)$.
    Verify $\pi P = \pi$ entry by entry:
    \begin{itemize}
        \item Vertex $1$: $\pi_2 \cdot \frac{1}{2} + \pi_3 \cdot \frac{1}{3} = \frac{1}{8} + \frac{1}{8} = \frac{1}{4}$.
        \item Vertex $2$: $\pi_1 \cdot \frac{1}{2} + \pi_3 \cdot \frac{1}{3} = \frac{1}{8} + \frac{1}{8} = \frac{1}{4}$.
        \item Vertex $3$: $\pi_1 \cdot \frac{1}{2} + \pi_2 \cdot \frac{1}{2} + \pi_4 \cdot 1
                  = \frac{1}{8} + \frac{1}{8} + \frac{1}{8} = \frac{3}{8}$.
        \item Vertex $4$: $\pi_3 \cdot \frac{1}{3} = \frac{1}{8}$.
    \end{itemize}

    (c) Yes. The graph is connected, and it contains a triangle (an odd cycle), so it is not
    bipartite.

    (d) $p_0 = (0, 0, 0, 1)$.
    \begin{itemize}
        \item $p_1 = (0, 0, 1, 0)$: from $4$ the walk must go to $3$.
        \item $p_2 = \left(\frac{1}{3}, \frac{1}{3}, 0, \frac{1}{3}\right)$: from $3$ each
              neighbour has probability $\frac{1}{3}$.
        \item $p_3 = \left(\frac{1}{3} \cdot \frac{1}{2},\ \frac{1}{3} \cdot \frac{1}{2},\
                  \frac{1}{3} \cdot \frac{1}{2} + \frac{1}{3} \cdot \frac{1}{2} + \frac{1}{3} \cdot 1,\ 0\right)
                  = \left(\frac{1}{6}, \frac{1}{6}, \frac{2}{3}, 0\right)$.
    \end{itemize}
    Each $p_t$ sums to $1$, as it must (a row stochastic matrix preserves probability vectors).
\end{wsolution}
\end{exercisebox}

\begin{exercisebox}
\begin{pproblem}[A path on three vertices]\label{pp:path}
    Consider the path $1 - 2 - 3$. Find $\pi$. Does the walk from vertex $1$ converge?
\end{pproblem}
\tcblower
\begin{wsolution}
    \leavevmode
    \begin{enumerate}
        \item The degrees are $1, 2, 1$ and $m = 2$, so
              $\pi = \left(\frac{1}{4}, \frac{1}{2}, \frac{1}{4}\right)$. Check:
              $\pi P = \left(\frac{1}{2} \cdot \frac{1}{2},\ \frac{1}{4} + \frac{1}{4},\ \frac{1}{2} \cdot \frac{1}{2}\right) = \pi$.
        \item The graph is bipartite, with sides $\set{1, 3}$ and $\set{2}$.
        \item From $p_0 = (1, 0, 0)$: $p_1 = (0, 1, 0)$,
              $p_2 = \left(\frac{1}{2}, 0, \frac{1}{2}\right)$, $p_3 = (0, 1, 0)$, and so on.
        \item The distribution alternates forever between $(0, 1, 0)$ and
              $\left(\frac{1}{2}, 0, \frac{1}{2}\right)$. It does not converge, although $\pi$
              exists and satisfies $\pi P = \pi$.
    \end{enumerate}
    In terms of eigenvalues: $P$ has the eigenvalues $1$, $0$ and $-1$. The eigenvalue $-1$ is what a
    bipartite graph produces, and its contribution $(-1)^t$ never dies out.
\end{wsolution}
\end{exercisebox}

\begin{exercisebox}
\begin{pproblem}[Short questions]\label{pp:walkshort}
    Find the stationary distribution and say whether the walk converges for
    \begin{enumerate}[label=(\alph*)]
        \item a $4$-cycle and a $5$-cycle,
        \item any $d$-regular graph on $n$ vertices,
        \item a disconnected graph.
    \end{enumerate}
\end{pproblem}
\tcblower
\begin{wsolution}
    (a) In a $4$-cycle every degree is $2$, so $\pi$ is uniform,
    $\left(\frac{1}{4}, \frac{1}{4}, \frac{1}{4}, \frac{1}{4}\right)$. An even cycle is bipartite, so
    the walk from a single vertex does not converge. A $5$-cycle is an odd cycle, so the walk does
    converge, to the uniform distribution.

    (b) $2m = nd$, so $\pi_v = \sfrac{d}{(nd)} = \sfrac{1}{n}$: the stationary distribution is
    uniform.

    (c) Each component has its own stationary distribution, so $\pi$ is not unique and the limit
    depends on the start.
\end{wsolution}
\end{exercisebox}

\section{Tail Bounds in a New Setting}\label{sec:tail}

Pick the inequality by what is known about $X$ (Sep 1, Sep 3).
\begin{center}
    \begin{adjustbox}{max width=\textwidth}
    \begin{tabular}{lll}
        \toprule
        Inequality & What it needs                                                      & Bound                                                           \\ \midrule
        Markov     & $X \geq 0$ and $\expectation{X}$                                   & $\prob{X \geq b} \leq \sfrac{\expectation{X}}{b}$               \\
        Chebyshev  & $\expectation{X}$ and $\variance{X}$                               & $\prob{\abs{X - \expectation{X}} \geq b} \leq \sfrac{\variance{X}}{b^2}$ \\
        Chernoff   & $X$ is a sum of independent $0/1$ variables, $\mu = \expectation{X}$ & exponentially small in $\mu$ (three forms below)                \\ \bottomrule
    \end{tabular}
    \end{adjustbox}
\end{center}
\begin{gather*}
    \prob{X \geq (1+\delta)\mu} \leq e^{-\mu\delta^2/3},
    \qquad
    \prob{X < (1-\delta)\mu} \leq e^{-\mu\delta^2/2}
    \quad (0 < \delta < 1), \\
    \prob{X \geq R} \leq 2^{-R} \quad (R \geq 6\mu).
\end{gather*}
For a statement about all $n$ bins, slots or coupons, bound one of them and multiply by $n$ (union
bound).

\begin{exercisebox}
\begin{pproblem}[The same event, three bounds]\label{pp:three}
    A fair coin is tossed $n$ times and $X$ is the number of heads. Bound
    $\prob{X \geq \sfrac{3n}{4}}$ with each inequality. Evaluate for $n = 100$.
\end{pproblem}
\tcblower
\begin{wsolution}
    $\expectation{X} = \sfrac{n}{2}$ and $\variance{X} = n \cdot \frac{1}{2} \cdot \frac{1}{2} = \sfrac{n}{4}$.
    \begin{enumerate}
        \item Markov: $\prob{X \geq \sfrac{3n}{4}} \leq \frac{n/2}{3n/4} = \frac{2}{3}$.
        \item Chebyshev: $X \geq \sfrac{3n}{4}$ implies $\abs{X - \sfrac{n}{2}} \geq \sfrac{n}{4}$, so
              the probability is at most $\frac{n/4}{(n/4)^2} = \frac{4}{n}$.
        \item Chernoff: $\sfrac{3n}{4} = \left(1 + \frac{1}{2}\right)\mu$ with $\mu = \sfrac{n}{2}$ and
              $\delta = \frac{1}{2}$, so the probability is at most
              $e^{-(n/2)(1/4)/3} = e^{-n/24}$.
    \end{enumerate}
    For $n = 100$: $0.667$, $0.04$ and $e^{-4.17} \approx 0.016$. Markov does not improve with $n$,
    Chebyshev improves like $\sfrac{1}{n}$, and Chernoff improves exponentially.
\end{wsolution}
\end{exercisebox}

\begin{exercisebox}
\begin{pproblem}[$n$ balls into $n$ bins]\label{pp:balls}
    Throw $n$ balls independently and uniformly into $n$ bins, $n \geq 4$. Show that with probability
    at least $1 - \sfrac{1}{n^2}$ no bin receives $3 \log_2 n$ balls or more.
\end{pproblem}
\tcblower
\begin{wsolution}
    \leavevmode
    \begin{enumerate}
        \item Fix a bin and let $X$ be its number of balls. $X$ is a sum of $n$ independent
              indicators, each $1$ with probability $\sfrac{1}{n}$, so $\mu = 1$.
        \item Take $R = 3 \log_2 n$. For $n \geq 4$ we have $R \geq 6 = 6\mu$, so the third Chernoff
              form applies.
        \item $\prob{X \geq 3 \log_2 n} \leq 2^{-3 \log_2 n} = \sfrac{1}{n^3}$.
        \item Union bound over the $n$ bins:
              $\prob{\text{some bin has at least } 3 \log_2 n \text{ balls}} \leq n \cdot \frac{1}{n^3} = \frac{1}{n^2}$. \qedhere
    \end{enumerate}
\end{wsolution}
\end{exercisebox}

\begin{exercisebox}
\begin{pproblem}[Coupon collector with a different number of trials]\label{pp:coupon}
    There are $n$ coupons and each trial gives a uniformly random one. Bound the probability that
    some coupon is missing after
    \begin{enumerate}[label=(\alph*)]
        \item $2n \ln n$ trials,
        \item $n \ln n + cn$ trials.
        \item What happens after only $n$ trials?
    \end{enumerate}
\end{pproblem}
\tcblower
\begin{wsolution}
    A fixed coupon is missed by one trial with probability $1 - \sfrac{1}{n}$, so it is missing
    after $m$ trials with probability $\left(1 - \frac{1}{n}\right)^m \leq e^{-m/n}$.

    (a) $e^{-2 \ln n} = \sfrac{1}{n^2}$ per coupon, so at most $n \cdot \frac{1}{n^2} = \frac{1}{n}$
    by the union bound.

    (b) $e^{-\ln n - c} = \sfrac{e^{-c}}{n}$ per coupon, so at most $e^{-c}$ in total. For $c = 3$
    this is about $0.05$.

    (c) The expected number of missing coupons is
    $n\left(1 - \frac{1}{n}\right)^n \approx \sfrac{n}{e}$, by linearity of expectation. About
    $37\%$ of the coupons are still missing.
\end{wsolution}
\end{exercisebox}

\begin{exercisebox}
\begin{pproblem}[Boosting a Monte Carlo algorithm by majority vote]\label{pp:majority}
    An algorithm gives the correct yes/no answer with probability $\sfrac{2}{3}$. Run it $t$ times
    independently ($t$ odd) and output the majority answer. How large must $t$ be for an error of at
    most $\sfrac{1}{n}$?
\end{pproblem}
\tcblower
\begin{wsolution}
    \leavevmode
    \begin{enumerate}
        \item Let $X$ be the number of correct runs: a sum of $t$ independent indicators with
              $\mu = \sfrac{2t}{3}$.
        \item The majority is wrong only if $X < \sfrac{t}{2}$.
        \item Write $\sfrac{t}{2} = (1 - \delta)\mu$: $(1 - \delta) \cdot \frac{2t}{3} = \frac{t}{2}$
              gives $\delta = \frac{1}{4}$.
        \item Lower tail:
              $\prob{X < \sfrac{t}{2}} \leq e^{-\mu\delta^2/2} = e^{-(2t/3)(1/16)/2} = e^{-t/48}$.
        \item $e^{-t/48} \leq \sfrac{1}{n}$ when $t \geq 48 \ln n$.
    \end{enumerate}
    So $\bigO{\log n}$ repetitions bring the error from $\sfrac{1}{3}$ down to $\sfrac{1}{n}$.
\end{wsolution}
\end{exercisebox}

\section{Lower Bounds by Adversary Argument and by Yao's Principle}\label{sec:lower}

There are two templates. \textbf{Adversary} (Aug 25), for deterministic algorithms: answer the
queries so that two different inputs stay consistent, and the algorithm must be wrong on one.
\textbf{Yao} (Sep 1, Oct 1), for randomized algorithms: choose a distribution on inputs and show
that every deterministic algorithm does badly on average; then every randomized algorithm does
badly on some input.

\begin{exercisebox}
\begin{pproblem}[Membership needs $n$ queries]\label{pp:membership}
    The scribe shows that finding $x$ needs $n - 1$ queries when $x$ is promised to be in the array.
    Show that deciding whether $x$ is in the array at all needs $n$ queries.
\end{pproblem}
\tcblower
\begin{wsolution}
    \leavevmode
    \begin{enumerate}
        \item Fix a deterministic algorithm. The adversary answers every query with ``not $x$''.
        \item Suppose the algorithm stops after at most $n - 1$ queries. Then some slot $j$ was never
              queried.
        \item Two arrays agree with all the answers: one has $x$ in slot $j$, the other has no $x$
              anywhere.
        \item The algorithm saw the same answers on both, so it gives the same output on both. It is
              wrong on one.
        \item So a correct algorithm makes $n$ queries.
    \end{enumerate}
    The promise is what saved one query on Aug 25: with a promise, $n - 1$ answers of ``no'' leave
    only one possible slot.
\end{wsolution}
\end{exercisebox}

\begin{exercisebox}
\begin{pproblem}[Deutsch--Jozsa needs $N/2 + 1$ deterministic queries]\label{pp:dj-lb}
    Prove that a deterministic classical algorithm for the Deutsch--Jozsa problem needs
    $2^{n-1} + 1$ queries.
\end{pproblem}
\tcblower
\begin{wsolution}
    \leavevmode
    \begin{enumerate}
        \item Let $N = 2^n$. The adversary answers $0$ to every query.
        \item Suppose the algorithm stops after at most $\sfrac{N}{2}$ queries. At least
              $\sfrac{N}{2}$ inputs are unqueried.
        \item Completion 1: $f = 0$ everywhere, which is constant.
        \item Completion 2: $f = 1$ on exactly $\sfrac{N}{2}$ of the unqueried inputs and $0$
              elsewhere, which is balanced.
        \item Both satisfy the promise and agree with every answer, so the algorithm is wrong on one. \qedhere
    \end{enumerate}
\end{wsolution}
\end{exercisebox}

\begin{exercisebox}
\begin{pproblem}[Randomized search needs about $n/2$ expected queries]\label{pp:yao-lv}
    Use Yao's principle to show that a Las Vegas algorithm for search makes at least
    $\sfrac{(n-1)}{2}$ expected queries on some input.
\end{pproblem}
\tcblower
\begin{wsolution}
    \leavevmode
    \begin{enumerate}
        \item Input distribution $\mu$: $x$ is in a uniformly random slot, each with probability
              $\sfrac{1}{n}$.
        \item Take a deterministic algorithm that is always correct. As long as every answer is
              ``no'', its queries follow one fixed order $q_1, q_2, \dots$
        \item It cannot stop while two or more slots are unqueried and $x$ has not been seen (the
              adversary argument). So if $x$ is in slot $q_i$ with $i \leq n - 1$, it makes at least
              $i$ queries, and if $x$ is in the last slot it makes $n - 1$.
        \item The expected number of queries under $\mu$ is at least
              \[\frac{1}{n} \left(\sum_{i=1}^{n-1} i + (n-1)\right) = \frac{(n-1)(n+2)}{2n} \geq \frac{n-1}{2}.\]
        \item By Yao's principle the expected cost of a randomized algorithm on its worst input is
              at least the expected cost of the best deterministic algorithm under $\mu$. So it is
              at least $\sfrac{(n-1)}{2} = \bigOm{n}$. \qedhere
    \end{enumerate}
\end{wsolution}
\end{exercisebox}

\begin{exercisebox}
\begin{pproblem}[The same for Monte Carlo algorithms]\label{pp:yao-mc}
    Show that a randomized search algorithm with success probability at least $\sfrac{2}{3}$ on
    every input needs at least $\sfrac{2n}{3} - 1$ queries.
\end{pproblem}
\tcblower
\begin{wsolution}
    \leavevmode
    \begin{enumerate}
        \item Use the same distribution $\mu$. Take a deterministic algorithm with $T$ queries, in
              the fixed order $q_1, \dots, q_T$ while the answers are ``no''.
        \item It finds $x$ if $x$ is in one of these $T$ slots: probability $\sfrac{T}{n}$.
        \item Otherwise it has seen $T$ answers of ``no'' and outputs one fixed slot. That slot is
              the right one with probability at most $\sfrac{1}{n}$ overall.
        \item So it succeeds with probability at most $\sfrac{(T+1)}{n}$ under $\mu$.
        \item If $T + 1 < \sfrac{2n}{3}$, every deterministic algorithm with $T$ queries errs with
              probability more than $\sfrac{1}{3}$ under $\mu$. By Yao's principle (the form proved
              on Oct 1) every randomized algorithm with $T$ queries errs with probability more than
              $\sfrac{1}{3}$ on some input. \qedhere
    \end{enumerate}
\end{wsolution}
\end{exercisebox}

\begin{exercisebox}
\begin{pproblem}[Simon's problem, in outline]\label{pp:simon-lb}
    State the hard distribution for Simon's problem and the key estimate. Evaluate it for $n = 10$
    and $16$ queries.
\end{pproblem}
\tcblower
\begin{wsolution}
    \leavevmode
    \begin{enumerate}
        \item Distribution: $s$ is uniform among the $N - 1$ nonzero strings, and $f$ is a uniformly
              random $2$-to-$1$ function with period $s$.
        \item An algorithm learns $s$ only from a collision, two queries $x$ and $x \oplus s$.
        \item With $k$ distinct queries the probability of a collision is at most
              $\sfrac{k(k-1)}{N}$, since there are $\sfrac{k(k-1)}{2}$ pairs and each pair rules in
              one value of $s$ out of at least $\sfrac{N}{2}$.
        \item Without a collision, $s$ is still uniform over at least $\sfrac{N}{2}$ strings, so a
              guess is right with probability at most $\sfrac{2}{N}$.
        \item The success probability is at most $\sfrac{(k(k-1) + 2)}{N}$.
    \end{enumerate}
    For $n = 10$, $N = 1024$ and $k = 16 = \sfrac{\sqrt{N}}{2}$: at most
    $\sfrac{(240 + 2)}{1024} \approx 0.24$. So $\sfrac{\sqrt{N}}{2}$ queries are not enough for
    success probability $\sfrac{2}{3}$, and the classical complexity is $\bigOm{\sqrt{2^n}}$.
\end{wsolution}
\end{exercisebox}

\section{Leader Election with a Changed Parameter}\label{sec:leader}

The algorithm (Sep 8) has two parameters: each node becomes a candidate with probability $p$, and
each candidate contacts $k$ random referees. In class $p = \sfrac{2 \ln n}{n}$ and
$k = \ceil{\sqrt{n \ln n}}$. The whole analysis is four estimates, and a changed parameter just
changes the numbers in them. Here $R_u$ is the referee set of node $u$.
\begin{gather*}
    \prob{\text{no candidate}} = (1-p)^n \leq e^{-pn},
    \qquad
    \prob{R_u \cap R_v = \emptyset} \leq \left(1 - \frac{k}{n}\right)^{k} \leq e^{-k^2/n}, \\
    \prob{\text{more than one leader}} \leq \binom{n}{2}\, p^2\, e^{-k^2/n},
    \qquad
    \expectation{\text{messages}} \leq 2k \cdot np.
\end{gather*}
\begin{center}
    \begin{adjustbox}{max width=\textwidth}
    \begin{tabular}{llll}
        \toprule
        Parameters                                                      & No leader                 & More than one leader                         & Expected messages                         \\ \midrule
        $p = \sfrac{2 \ln n}{n}$, $k = \ceil{\sqrt{n \ln n}}$ (class)   & $\leq \sfrac{1}{n^2}$     & $\leq \sfrac{2 \ln^2 n}{n}$                  & $\leq 4k \ln n$                           \\
        $p = \sfrac{3 \ln n}{n}$, same $k$                              & $\leq \sfrac{1}{n^3}$     & $\leq \sfrac{9 \ln^2 n}{(2n)}$               & $\leq 6k \ln n$                           \\
        same $p$, $k = \ceil{2\sqrt{n \ln n}}$                          & $\leq \sfrac{1}{n^2}$     & $\leq \sfrac{2 \ln^2 n}{n^4}$                & $\leq 4k \ln n$, with $k$ twice as large  \\
        same $p$, $k = \ceil{\sqrt{n}}$                                 & $\leq \sfrac{1}{n^2}$     & bound is $\sfrac{2 \ln^2 n}{e}$, useless     & $\leq 4k \ln n$                           \\
        $p = \sfrac{1}{n}$, same $k$                                    & about $\sfrac{1}{e}$      & $\leq \sfrac{1}{(2n)}$                       & $\leq 2k$                                 \\ \bottomrule
    \end{tabular}
    \end{adjustbox}
\end{center}

\begin{exercisebox}
\begin{pproblem}[$p = 3 \ln n / n$]\label{pp:leader3}
    Redo the analysis with candidate probability $p = \sfrac{3 \ln n}{n}$ and the same $k$.
\end{pproblem}
\tcblower
\begin{wsolution}
    \leavevmode
    \begin{enumerate}
        \item At least one leader: the candidate with the smallest id always wins at all its
              referees, so there is no leader only if there is no candidate.
              $\prob{\text{no candidate}} = (1-p)^n \leq e^{-pn} = e^{-3 \ln n} = \sfrac{1}{n^3}$.
        \item Two candidates with a common referee cannot both be leaders. Two referee sets are
              disjoint with probability at most $e^{-k^2/n} \leq e^{-\ln n} = \sfrac{1}{n}$, as
              before.
        \item A fixed pair is bad (both are candidates, with no common referee) with probability at
              most $\sfrac{p^2}{n}$. The expected number of bad pairs is at most
              \[\binom{n}{2} \frac{p^2}{n} \leq \frac{n^2}{2} \cdot \frac{9 \ln^2 n}{n^2} \cdot \frac{1}{n} = \frac{9 \ln^2 n}{2n}.\]
        \item By Markov,
              $\prob{\text{more than one leader}} \leq \prob{\text{at least one bad pair}} \leq \frac{9 \ln^2 n}{2n}$.
        \item Messages: at most $2k$ per candidate and $3 \ln n$ candidates in expectation, so at
              most $6k \ln n = \bigO{\sqrt{n} \log^{3/2} n}$.
    \end{enumerate}
    More candidates make ``no leader'' less likely, and ``two leaders'' and the message count
    slightly larger. The asymptotics do not change.
\end{wsolution}
\end{exercisebox}

\begin{exercisebox}
\begin{pproblem}[Only $\sqrt{n}$ referees]\label{pp:leader-sqrt}
    What breaks if $k = \ceil{\sqrt{n}}$, without the logarithm?
\end{pproblem}
\tcblower
\begin{wsolution}
    \leavevmode
    \begin{enumerate}
        \item Now $\sfrac{k^2}{n} \approx 1$, so two referee sets are disjoint with probability about
              $e^{-1} \approx 0.37$ instead of $\sfrac{1}{n}$.
        \item The bound on the expected number of bad pairs becomes
              $\frac{n^2}{2} \cdot \frac{4 \ln^2 n}{n^2} \cdot e^{-1} = \frac{2 \ln^2 n}{e}$, which is
              larger than $1$. Markov says nothing.
        \item The algorithm itself fails, not only the proof. Take the two candidates with the
              smallest ids. With probability about $0.37$ they share no referee. Then, unless the
              second one happened to pick the first as a referee, none of its referees hears a
              smaller id, and both declare themselves leader.
    \end{enumerate}
    This is the birthday paradox: two random sets of size $k$ out of $n$ meet reliably only when
    $k^2$ is larger than $n$ by a log factor.
\end{wsolution}
\end{exercisebox}

\begin{exercisebox}
\begin{pproblem}[Too few candidates]\label{pp:leader-few}
    What happens with $p = \sfrac{1}{n}$?
\end{pproblem}
\tcblower
\begin{wsolution}
    The expected number of candidates is $1$, and
    $\prob{\text{no candidate}} = \left(1 - \frac{1}{n}\right)^n$, which tends to
    $\sfrac{1}{e} \approx 0.37$. So with constant probability nobody is elected. The factor $\ln n$
    in $p$ is what pushes this probability down to $\sfrac{1}{n^2}$.
\end{wsolution}
\end{exercisebox}

\begin{exercisebox}
\begin{pproblem}[Numbers for $n = 1{,}000{,}000$]\label{pp:leader-num}
    Evaluate the class algorithm for $n = 10^6$.
\end{pproblem}
\tcblower
\begin{wsolution}
    $\ln n \approx 13.82$.
    \begin{enumerate}
        \item Expected candidates: $np = 2 \ln n \approx 27.6$.
        \item Referees per candidate: $k = \ceil{\sqrt{10^6 \cdot 13.82}} = 3717$.
        \item Expected messages: at most $4k \ln n \approx 205{,}000$. The trivial algorithm, where
              everyone sends its id to everyone, uses about $n^2 = 10^{12}$.
        \item Failure probability: at most
              $\frac{1}{n^2} + \frac{2 \ln^2 n}{n} \approx 10^{-12} + 3.8 \times 10^{-4}$.
        \item Rounds: $3$. \qedhere
    \end{enumerate}
\end{wsolution}
\end{exercisebox}

\section{Deterministic, Randomized and Quantum Query Counts}\label{sec:counts}

All entries are from the scribes. $N = 2^n$ is the number of inputs.
\begin{center}
    \begin{adjustbox}{max width=\textwidth}
    \begin{tabular}{lllll}
        \toprule
        Problem                               & Deterministic         & Randomized                       & Quantum             & Lecture              \\ \midrule
        Search, one solution among $N$        & $N - 1$               & $\bigOm{N}$                      & $\bigO{\sqrt{N}}$   & Aug 25, Sep 1, Oct 6 \\
        Deutsch                               & $2$                   & $2$                              & $1$                 & Sep 24               \\
        Deutsch--Jozsa                        & $\sfrac{N}{2} + 1$    & $\bigO{1}$ for constant error    & $1$, no error       & Sep 24               \\
        Bernstein--Vazirani                   & $n$                   & $n$                              & $1$                 & Sep 29               \\
        Simon                                 & $\bigTh{\sqrt{N}}$    & $\bigTh{\sqrt{N}}$               & $\bigO{n}$          & Sep 29, Oct 1        \\
        Collision, random function into $[n]$ & ---                   & $\bigTh{\sqrt{n}}$               & later in the course & Sep 3                \\ \bottomrule
    \end{tabular}
    \end{adjustbox}
\end{center}

\begin{exercisebox}
\begin{pproblem}[Is Deutsch--Jozsa an exponential speedup?]\label{pp:dj-speedup}
    The quantum algorithm uses $1$ query and a deterministic one needs $\sfrac{N}{2} + 1$. Is this
    an exponential quantum advantage?
\end{pproblem}
\tcblower
\begin{wsolution}
    Only against deterministic (or zero-error) algorithms. A randomized algorithm queries $\alpha$
    random inputs and answers ``balanced'' if two answers differ. It is always right on a constant
    $f$. On a balanced $f$ all $\alpha$ answers agree with probability
    $2 \cdot \left(\frac{1}{2}\right)^{\alpha} = \left(\frac{1}{2}\right)^{\alpha - 1}$. With
    $\alpha = 11$ queries the error is below $\sfrac{1}{1000}$, for every $n$. So against
    bounded-error randomized algorithms the gap is $1$ query against a constant.
\end{wsolution}
\end{exercisebox}

\begin{exercisebox}
\begin{pproblem}[Why Simon is different]\label{pp:simon-sep}
    Explain why Simon's problem gives a real exponential separation.
\end{pproblem}
\tcblower
\begin{wsolution}
    The Oct 1 lower bound is for randomized algorithms with bounded error: success probability
    $\sfrac{2}{3}$ needs more than $\sfrac{\sqrt{N}}{2}$ queries. The quantum algorithm needs
    $\bigO{n}$ queries. So the gap is $\sqrt{2^n}$ against $n$ even when the classical algorithm may
    use randomness and may err. A quantum advantage has to be shown against randomized algorithms,
    because they are the strongest classical model (Aug 25).
\end{wsolution}
\end{exercisebox}

\needspace{7cm}
\begin{exercisebox}
\begin{pproblem}[The numbers for $n = 20$]\label{pp:n20}
    Fill in the query counts for $n = 20$, $N = 1{,}048{,}576$.
\end{pproblem}
\tcblower
\begin{wsolution}
    The classical count for Simon lies between the lower bound of Oct 1 and the deterministic
    algorithm of Sep 29.
    \begin{center}
        \begin{tabular}{lp{4.4cm}p{4.4cm}}
            \toprule
            Problem                   & Classical                                                                         & Quantum                                                         \\ \midrule
            Deutsch--Jozsa, no error  & $524{,}289$                                                                       & $1$                                                             \\
            Bernstein--Vazirani       & $20$                                                                              & $1$                                                             \\
            Simon                     & more than $512$ ($\sfrac{\sqrt{N}}{2}$); about $1024$ suffice up to a constant    & $21$ per attempt: $19$ quantum and $2$ classical                \\
            Search, one solution      & $1{,}048{,}575$ deterministic                                                     & $804$ Grover iterations: $\floor{\frac{\pi}{4} \cdot 1024}$     \\ \bottomrule
        \end{tabular}
    \end{center}
\end{wsolution}
\end{exercisebox}

\begin{exercisebox}
\begin{pproblem}[One-line reasons]\label{pp:reasons}
    Give a one-line reason for each statement.
    \begin{enumerate}[label=(\alph*)]
        \item Search is only a quadratic speedup.
        \item Bernstein--Vazirani is $n$ against $1$.
        \item Simon and collisions have the same square root.
    \end{enumerate}
\end{pproblem}
\tcblower
\begin{wsolution}
    (a) Grover uses $\bigO{\sqrt{N}}$ queries and the scribe notes that $\bigOm{\sqrt{N}}$ quantum
    queries are necessary, so there is no exponential gain for unstructured search.

    (b) Each classical answer is one bit and $s$ has $n$ bits. One quantum query returns all $n$
    bits because the Hadamard transform turns the sign pattern $(-1)^{s \cdot x}$ into $\ket{s}$.

    (c) Classically $s$ is found only by a collision, and $k$ queries make about $\sfrac{k^2}{2}$
    pairs, so $k$ must reach about $\sqrt{N}$. This is the birthday paradox of Sep 3.
\end{wsolution}
\end{exercisebox}

\section{The Chernoff Bound: Proof}\label{sec:chernoff}

The scribe states the Chernoff bounds (Sep 1) without proof. The proof is three ideas:
exponentiate, apply Markov, and use independence to turn the expectation of a product into a
product of expectations. This is the standard proof; compare with slides 3--7 of the Chernoff deck.

\begin{theorem}[Chernoff Bounds]\label{thm:chernoff-pp}
    Let $X = X_1 + \dots + X_n$ with the $X_i$ independent, $X_i = 1$ with probability $p_i$ and $0$
    otherwise, and let $\mu = \expectation{X} = p_1 + \dots + p_n > 0$. Then
    \begin{enumerate}
        \item $\prob{X \geq (1+\delta)\mu} \leq e^{-\mu\delta^2/3}$ for $0 < \delta < 1$,
        \item $\prob{X < (1-\delta)\mu} \leq e^{-\mu\delta^2/2}$ for $0 < \delta < 1$,
        \item $\prob{X \geq R} \leq 2^{-R}$ for $R \geq 6\mu$.
    \end{enumerate}
\end{theorem}

\begin{proof}
    \emph{Step 1: Markov on $e^{tX}$.} For $t > 0$ the function $e^{tx}$ is increasing, so
    $X \geq a$ is the same event as $e^{tX} \geq e^{ta}$. Since $e^{tX} > 0$, Markov's inequality
    gives
    \[\prob{X \geq a} = \prob{e^{tX} \geq e^{ta}} \leq \frac{\expectation{e^{tX}}}{e^{ta}}.\]

    \emph{Step 2: independence.} $e^{tX}$ is the product of the $e^{tX_i}$, and these are
    independent, so the expectation factors. Each factor is
    $\expectation{e^{tX_i}} = p_i e^t + (1 - p_i) = 1 + p_i(e^t - 1)$. Using $1 + x \leq e^x$,
    \[\expectation{e^{tX}} = \prod_{i=1}^{n} \left(1 + p_i(e^t - 1)\right)
        \leq \prod_{i=1}^{n} e^{p_i(e^t - 1)} = e^{\mu(e^t - 1)}.\]
    This holds for every real $t$, positive or negative.

    \emph{Step 3: choose $t$.} Put $a = (1+\delta)\mu$ and $t = \ln(1+\delta)$, which is positive
    for $\delta > 0$. Then $e^t - 1 = \delta$ and $e^{ta} = (1+\delta)^{(1+\delta)\mu}$, so for every
    $\delta > 0$
    \begin{equation}\label{eq:chernoff-general}
        \prob{X \geq (1+\delta)\mu} \leq \left(\frac{e^{\delta}}{(1+\delta)^{1+\delta}}\right)^{\mu}.
    \end{equation}
    This is the general form. The three forms of the theorem are simplifications of it.

    \emph{Step 4: the upper tail for $0 < \delta < 1$.} We need
    $\delta - (1+\delta)\ln(1+\delta) \leq -\sfrac{\delta^2}{3}$. Let
    \[g(\delta) = \delta - (1+\delta)\ln(1+\delta) + \frac{\delta^2}{3},
        \qquad
        g'(\delta) = -\ln(1+\delta) + \frac{2\delta}{3},
        \qquad
        g''(\delta) = -\frac{1}{1+\delta} + \frac{2}{3}.\]
    We have $g(0) = 0$ and $g'(0) = 0$. Also $g''$ is negative for $\delta < \frac{1}{2}$ and
    positive for $\delta > \frac{1}{2}$, so $g'$ first decreases and then increases. Since
    $g'(1) = \frac{2}{3} - \ln 2 \approx -0.026 < 0$, the function $g'$ starts at $0$, goes down, and
    has not come back up to $0$ by $\delta = 1$: $g' \leq 0$ on $[0, 1]$. Hence $g$ is decreasing
    on $[0, 1]$ and $g(\delta) \leq g(0) = 0$. Therefore
    $\prob{X \geq (1+\delta)\mu} \leq e^{-\mu\delta^2/3}$.

    \emph{Step 5: the form for $R \geq 6\mu$.} Write $R = (1+\delta)\mu$, so $1 + \delta \geq 6$.
    In \cref{eq:chernoff-general}, bound $e^{\delta}$ by $e^{1+\delta}$:
    \[\prob{X \geq R} \leq \left(\frac{e}{1+\delta}\right)^{(1+\delta)\mu}
        = \left(\frac{e}{1+\delta}\right)^{R} \leq \left(\frac{e}{6}\right)^{R} \leq 2^{-R},\]
    because $\sfrac{e}{6} \approx 0.453 < \frac{1}{2}$.

    \emph{Step 6: the lower tail for $0 < \delta < 1$.} Now take $t < 0$. Then $e^{tx}$ is
    decreasing, so $X \leq a$ is the same event as $e^{tX} \geq e^{ta}$, and Steps 1 and 2 give
    $\prob{X \leq a} \leq e^{\mu(e^t - 1) - ta}$. Put $a = (1-\delta)\mu$ and
    $t = \ln(1-\delta) < 0$:
    \[\prob{X \leq (1-\delta)\mu} \leq \left(\frac{e^{-\delta}}{(1-\delta)^{1-\delta}}\right)^{\mu}.\]
    We need $-\delta - (1-\delta)\ln(1-\delta) \leq -\sfrac{\delta^2}{2}$. Let
    $h(\delta) = -\delta - (1-\delta)\ln(1-\delta) + \sfrac{\delta^2}{2}$. Then $h(0) = 0$ and
    \[h'(\delta) = -1 + \ln(1-\delta) + 1 + \delta = \ln(1-\delta) + \delta \leq 0,\]
    since $\ln(1-\delta) \leq -\delta$. So $h \leq 0$ on $[0, 1)$, and
    $\prob{X < (1-\delta)\mu} \leq \prob{X \leq (1-\delta)\mu} \leq e^{-\mu\delta^2/2}$.
\end{proof}

\paragraph{What to remember.}
\begin{itemize}
    \item Independence is used exactly once, in Step 2. Without it the proof fails, and only Markov
          or Chebyshev are available.
    \item For $\delta \geq 1$ use the general form \cref{eq:chernoff-general} or the $2^{-R}$ form.
          The $e^{-\mu\delta^2/3}$ form is only proved for $\delta < 1$.
\end{itemize}

\section{Max Load, Upper Tail: $3n \ln n$ Queries, at Most $9 \ln n$ per Slot}\label{sec:maxload}

The Sep 1 scribe proves the lower tail (no slot is accessed too rarely, with $24 n \ln n$ queries).
The version in the professor's own notes is the upper tail with $3n \ln n$ queries, the same number
that makes every slot accessed at least once.

\begin{claim}[Max Load, Upper Tail]\label{clm:maxload-pp}
    Random sampling makes $3n \ln n$ queries into an array with $n \geq 2$ slots. Then with
    probability at least $1 - \sfrac{1}{n^2}$ no slot is accessed $9 \ln n$ times or more.
\end{claim}

\begin{proof}
    Fix a slot $i$ and let $Y$ be the number of queries that access it. $Y$ is a sum of $3n \ln n$
    independent indicators, each $1$ with probability $\sfrac{1}{n}$, so
    $\mu = \expectation{Y} = 3 \ln n$. The threshold is $9 \ln n = 3\mu = (1+\delta)\mu$ with
    $\delta = 2$. Since $\delta = 2$ is not below $1$, we use the general form
    \cref{eq:chernoff-general}:
    \[\prob{Y \geq 9 \ln n} \leq \left(\frac{e^{2}}{3^{3}}\right)^{3 \ln n}
        = n^{3(2 - \ln 27)} = n^{-3.88\ldots} \leq \frac{1}{n^{3}}.\]
    Here $\ln 27 \approx 3.296$, so the exponent is $3 \cdot (2 - 3.296) \approx -3.89$. By the
    union bound over the $n$ slots,
    \[\prob{\text{some slot is accessed at least } 9 \ln n \text{ times}} \leq n \cdot \frac{1}{n^3} = \frac{1}{n^2}. \qedhere\]
\end{proof}

This is stronger than the $1 - \sfrac{1}{n}$ in the professor's notes.

\paragraph{With only the three forms in the scribe.}
If only the stated forms may be used, take $R = 6\mu = 18 \ln n$ in the third one:
$\prob{Y \geq 18 \ln n} \leq 2^{-18 \ln n} = n^{-18 \ln 2} \approx n^{-12.5}$. The union bound gives
a failure probability of at most $n^{-11}$. The constant is $18$ instead of $9$, but no new formula
is needed.

\paragraph{Putting both tails together.}
By the Sep 1 theorem, $3n \ln n$ queries access every slot at least once with probability at least
$1 - \sfrac{1}{n^2}$. With \cref{clm:maxload-pp} and a union bound, with probability at least
$1 - \sfrac{2}{n^2}$ every slot is accessed at least once and fewer than $9 \ln n$ times.

\paragraph{Numbers for $n = 1000$.}
$\ln n \approx 6.91$, so the algorithm makes $3n \ln n \approx 20{,}700$ queries. Each slot is
accessed about $20.7$ times on average. The cap $9 \ln n$ is about $62$. The probability that some
slot reaches the cap is at most $1000^{-2.89} \approx 2 \times 10^{-9}$.

\section{Birthday Paradox and Collisions}\label{sec:birthday}

Everything here is from Sep 3. There are $m$ people (queries) and $n$ days (values), and $X$ is the
number of colliding pairs.
\begin{gather*}
    \expectation{X} = \binom{m}{2} \frac{1}{n} = \frac{m(m-1)}{2n},
    \qquad
    \prob{\text{no collision}} = \prod_{i=0}^{m-1} \left(1 - \frac{i}{n}\right) \leq e^{-m(m-1)/(2n)}, \\
    \prob{X = 0} \leq \frac{\variance{X}}{\expectation{X}^2} \leq \frac{1}{\expectation{X}}.
\end{gather*}
The first is linearity of expectation, the second is the chain rule with $1 - x \leq e^{-x}$, and
the third is Chebyshev with $\variance{X} \leq \expectation{X}$.

\begin{exercisebox}
\begin{pproblem}[The numbers for $n = 365$]\label{pp:bday365}
    \leavevmode
    \begin{enumerate}[label=(\alph*)]
        \item How many people make the expected number of collisions at least $1$?
        \item How many make the probability of a collision at least $\frac{1}{2}$, by the chain-rule
              bound?
        \item How many does the Chebyshev proof need to guarantee probability $\frac{1}{2}$?
    \end{enumerate}
\end{pproblem}
\tcblower
\begin{wsolution}
    (a) We need $m(m-1) \geq 2n = 730$. Since $27 \cdot 26 = 702$ and $28 \cdot 27 = 756$, the
    answer is $m = 28$, with $\expectation{X} = \sfrac{756}{730} \approx 1.04$.

    (b) We need $e^{-m(m-1)/730} \leq \frac{1}{2}$, that is $m(m-1) \geq 730 \ln 2 \approx 505.997$.
    Since $22 \cdot 21 = 462$ and $23 \cdot 22 = 506$, the answer is $m = 23$. The exact product for
    $m = 23$ is $0.4927$, so a collision has probability $0.507$.

    (c) Chebyshev gives $\prob{X = 0} \leq \sfrac{1}{\expectation{X}}$, so we need
    $\expectation{X} \geq 2$, that is $m(m-1) \geq 1460$. Since $38 \cdot 37 = 1406$ and
    $39 \cdot 38 = 1482$, the answer is $m = 39$. The condition $m \geq 2\sqrt{n} + 1 \approx 39.2$
    of the exercise gives $m = 40$.

    The three answers differ ($28$, $23$, $39$) but all are about $\sqrt{n}$. Part (a) alone proves
    nothing about the probability: $\expectation{X} \geq 1$ does not imply that
    $\prob{X \geq 1}$ is large. That step needs more than the expectation: either the direct
    product of part (b), or the variance, which is what the exercise uses with Chebyshev.
\end{wsolution}
\end{exercisebox}

\begin{exercisebox}
\begin{pproblem}[Why Chebyshev and not Chernoff?]\label{pp:why-cheb}
    $X$ is a sum of indicators $X_{ij}$. Why does the birthday proof not use the Chernoff bound?
    Why is Markov not enough?
\end{pproblem}
\tcblower
\begin{wsolution}
    \leavevmode
    \begin{enumerate}
        \item Chernoff needs mutually independent indicators. The $X_{ij}$ are not: if persons $1$
              and $2$ share a birthday and persons $2$ and $3$ share a birthday, then persons $1$
              and $3$ do too. So $\prob{X_{12} = X_{23} = X_{13} = 1} = \sfrac{1}{n^2}$, not
              $\sfrac{1}{n^3}$.
        \item They are pairwise independent: for two different pairs,
              $\prob{X_{ij} = X_{kl} = 1} = \sfrac{1}{n^2} = \prob{X_{ij} = 1} \cdot \prob{X_{kl} = 1}$,
              whether the pairs share a person or not.
        \item Pairwise independence is exactly what the variance needs. The cross terms
              $\expectation{X_{ij} X_{kl}} - \expectation{X_{ij}}\expectation{X_{kl}}$ are all zero,
              so $\variance{X}$ is the sum of the $\variance{X_{ij}}$, which is at most
              $\expectation{X}$. Chebyshev then applies.
        \item Markov gives $\prob{X \geq 1} \leq \expectation{X}$. That is an upper bound on the
              probability of a collision, and we want a lower bound.
    \end{enumerate}
    The rule: Markov needs only $X \geq 0$, Chebyshev needs the variance (pairwise independence is
    enough), and Chernoff needs mutual independence.
\end{wsolution}
\end{exercisebox}

\begin{exercisebox}
\begin{pproblem}[The collision lower bound with numbers]\label{pp:collision-num}
    Let $f$ be a random function into $n = 1{,}000{,}000$ values.
    \begin{enumerate}[label=(\alph*)]
        \item Bound the probability that $500$ queries find a collision.
        \item In general, bound it for $k = \sfrac{\sqrt{n}}{2}$ queries.
        \item How many queries are needed for probability $\frac{1}{2}$, and how many are enough?
    \end{enumerate}
\end{pproblem}
\tcblower
\begin{wsolution}
    By the Sep 3 exercise, $k$ queries find a collision with probability at most
    $\sfrac{k(k-1)}{(2n)}$. The reason is that query number $s$ hits one of the $s - 1$ earlier
    answers with probability at most $\sfrac{(s-1)}{n}$, and the union bound adds these up.

    (a) $\frac{500 \cdot 499}{2{,}000{,}000} \approx 0.125$.

    (b) $\frac{k(k-1)}{2n} < \frac{k^2}{2n} = \frac{n/4}{2n} = \frac{1}{8}$, for every $n$.

    (c) Needed: $\frac{k(k-1)}{2n} \geq \frac{1}{2}$ means $k(k-1) \geq 1{,}000{,}000$. Since
    $1000 \cdot 999 = 999{,}000$ is too small, $k \geq 1001$. Enough: by the birthday exercise,
    $m \geq 2\sqrt{n} + 1 = 2001$ random queries collide with probability at least $\frac{1}{2}$.

    So the answer lies between $1001$ and $2001$ queries: $\bigTh{\sqrt{n}}$. These are two
    separate proofs. The upper bound analyses one algorithm (random queries, Chebyshev). The lower
    bound holds for every algorithm, adaptive or randomized (union bound).
\end{wsolution}
\end{exercisebox}

\section{Expectation, Markov and the Weak Law}\label{sec:expect}

Three short problems that each use one tool: linearity of expectation (Aug 27), the weak law from
Chebyshev (Sep 8), and Markov (Sep 3).

\begin{exercisebox}
\begin{pproblem}[Empty bins]\label{pp:empty}
    Throw $m$ balls independently and uniformly into $n$ bins.
    \begin{enumerate}[label=(\alph*)]
        \item Find the expected number of empty bins.
        \item Evaluate it for $m = n$, $m = n \ln n$ and $m = 2n \ln n$.
        \item For $m = 2n \ln n$, bound the probability that some bin is empty.
    \end{enumerate}
\end{pproblem}
\tcblower
\begin{wsolution}
    Let $Y_i = 1$ if bin $i$ is empty and $0$ otherwise, and let $Y = Y_1 + \dots + Y_n$ be the
    number of empty bins. One ball misses bin $i$ with probability $1 - \sfrac{1}{n}$, and the balls
    are independent, so
    $\expectation{Y_i} = \prob{Y_i = 1} = \left(1 - \frac{1}{n}\right)^m$.

    (a) By linearity of expectation,
    $\expectation{Y} = n\left(1 - \frac{1}{n}\right)^m \leq n e^{-m/n}$, using $1 - x \leq e^{-x}$.

    (b) For $m = n$: about $\sfrac{n}{e}$, so $37\%$ of the bins are empty (for $n = 100$ the exact
    value is $36.6$). For $m = n \ln n$: at most $n e^{-\ln n} = 1$. For $m = 2n \ln n$: at most
    $n \cdot n^{-2} = \sfrac{1}{n}$.

    (c) $Y \geq 0$, so Markov gives $\prob{Y \geq 1} \leq \expectation{Y} \leq \sfrac{1}{n}$.

    The $Y_i$ are not independent (if every other bin is empty, bin $i$ is not), and linearity does
    not need them to be. Part (c) is the coupon collector bound again: Markov on the number of empty
    bins gives the same number as the union bound over the bins.
\end{wsolution}
\end{exercisebox}

\begin{exercisebox}
\begin{pproblem}[Estimating the bias of a coin]\label{pp:coin}
    A coin shows heads with unknown probability $p$. We toss it $n$ times and estimate $p$ by the
    fraction of heads. How large must $n$ be so that the estimate is within $\eps$ of $p$ with
    probability at least $1 - \delta$? Evaluate for $\eps = \delta = 0.05$.
\end{pproblem}
\tcblower
\begin{wsolution}
    Let $X_i = 1$ if toss $i$ is heads. Then $\expectation{X_i} = p$ and
    $\variance{X_i} = p(1-p) \leq \frac{1}{4}$, because $p(1-p)$ is largest at $p = \frac{1}{2}$. The
    estimate is the average $\overline{X}_n$ of $X_1, \dots, X_n$. By the weak law (the Sep 8
    exercise, which is Chebyshev applied to the average),
    \[\prob{\abs{\overline{X}_n - p} \geq \eps} \leq \frac{\sigma^2}{n\eps^2}
        = \frac{p(1-p)}{n\eps^2} \leq \frac{1}{4n\eps^2}.\]
    This is at most $\delta$ when $n \geq \frac{1}{4\eps^2\delta}$. For $\eps = \delta = 0.05$:
    $n \geq \frac{1}{4 \cdot 0.0025 \cdot 0.05} = 2000$ tosses.

    The bound does not need $p$, which is the point: $p$ is what we are trying to learn. The weak
    law only needs pairwise independent tosses. The cost of using Chebyshev is the factor
    $\sfrac{1}{\delta}$: for $\delta = 10^{-6}$ it asks for $10^8$ tosses. Independent tosses allow
    Chernoff, where $\delta$ enters only as $\ln(\sfrac{1}{\delta})$.
\end{wsolution}
\end{exercisebox}

\begin{exercisebox}
\begin{pproblem}[Las Vegas to Monte Carlo]\label{pp:lv-mc}
    A Las Vegas algorithm $\mathcal{A}$ is always correct and makes $T$ queries in expectation on
    every input.
    \begin{enumerate}[label=(\alph*)]
        \item Build an algorithm that never makes more than $2T$ queries and fails with probability
              at most $\frac{1}{2}$.
        \item Get the failure probability down to $\sfrac{1}{n}$.
        \item Check (a) on random sampling for search.
    \end{enumerate}
\end{pproblem}
\tcblower
\begin{wsolution}
    (a) Run $\mathcal{A}$ and stop it after $2T$ queries. If it has finished, output its answer,
    which is correct. If not, output ``fail''. Let $X$ be the number of queries that $\mathcal{A}$
    would make. $X \geq 0$ and $\expectation{X} \leq T$, so by Markov
    $\prob{X \geq 2T} \leq \frac{\expectation{X}}{2T} \leq \frac{1}{2}$. The new algorithm fails only
    in that case. Stopping after $cT$ queries gives failure at most $\sfrac{1}{c}$ in the same way.

    (b) Markov alone would need $c = n$, that is $nT$ queries. Instead repeat the $2T$-query version
    $\ceil{\log_2 n}$ times with fresh randomness and stop at the first run that finishes. All runs
    fail with probability at most $\left(\frac{1}{2}\right)^{\log_2 n} = \frac{1}{n}$. The total is
    at most $2T \ceil{\log_2 n}$ queries.

    (c) Random sampling with replacement finds $x$ with probability $\sfrac{1}{n}$ per query, so the
    number of queries is geometric with mean $T = n$. Markov says that $2n$ queries fail with
    probability at most $\frac{1}{2}$. The exact value is
    $\left(1 - \frac{1}{n}\right)^{2n} \leq e^{-2} \approx 0.135$.

    Markov is weak here, but it used nothing about $\mathcal{A}$ except its expected cost. Part (b)
    works because a Las Vegas run tells us whether it finished, so failures are detected and never
    output as answers.
\end{wsolution}
\end{exercisebox}

\section{Three More Quantum Questions}\label{sec:more}

\begin{exercisebox}
\begin{pproblem}[Why do we need $U_f$?]\label{pp:uf}
    Let $f \colon \set{0,1}^n \to \set{0,1}$.
    \begin{enumerate}[label=(\alph*)]
        \item Why can no quantum circuit map $\ket{x}$ to $\ket{f(x)}$?
        \item Show that $U_f \ket{x}\ket{y} = \ket{x}\ket{y \oplus f(x)}$ is unitary.
        \item How do we read $f(x)$ from it, and where does the sign $(-1)^{f(x)}$ come from?
    \end{enumerate}
\end{pproblem}
\tcblower
\begin{wsolution}
    (a) Every quantum gate is unitary, so it preserves norms and can be undone (Sep 15). The map
    $x \mapsto f(x)$ is not one-to-one as soon as two inputs have the same value, and with more than
    two inputs that always happens. In vectors: if $f(x) = f(x')$ with $x \neq x'$, a linear map $T$
    with $T\ket{x} = \ket{f(x)}$ sends $\ket{x} - \ket{x'}$ to $0$. That vector has norm $\sqrt{2}$,
    so $T$ does not preserve norms. This is the Sep 15 argument for XOR, where $00$ and $11$ both
    give $0$.

    (b) $U_f$ sends the basis state $\ket{x}\ket{y}$ to the basis state
    $\ket{x}\ket{y \oplus f(x)}$. Applying it twice gives $y \oplus f(x) \oplus f(x) = y$, so
    $U_f^2 = I$. So $U_f$ permutes the $2^{n+1}$ basis states, its matrix is a permutation matrix,
    and a permutation matrix is unitary (the Sep 15 argument for CNOT). The trick is that $x$ is
    kept: from $(x, y \oplus f(x))$ we can recover $y$, so nothing is lost.

    (c) With $y = 0$: $U_f \ket{x}\ket{0} = \ket{x}\ket{f(x)}$, so measuring the last qubit gives
    $f(x)$. With the last qubit in $\ket{-} = \frac{1}{\sqrt{2}}\left(\ket{0} - \ket{1}\right)$:
    \[U_f \ket{x}\ket{-} = \ket{x} \cdot \frac{1}{\sqrt{2}}\left(\ket{f(x)} - \ket{1 \oplus f(x)}\right).\]
    If $f(x) = 0$ this is $\ket{x}\ket{-}$, and if $f(x) = 1$ it is $-\ket{x}\ket{-}$. So
    $U_f \ket{x}\ket{-} = (-1)^{f(x)} \ket{x}\ket{-}$, the first fact of \cref{sec:run}.

    CNOT is the smallest example: it is $U_f$ for the one-bit function $f(x) = x$.
\end{wsolution}
\end{exercisebox}

\begin{exercisebox}
\begin{pproblem}[Deterministic Simon for $n = 4$]\label{pp:det-simon}
    Run the deterministic algorithm of the Sep 29 exercise for $n = 4$.
    \begin{enumerate}[label=(\alph*)]
        \item Which inputs does it query, and how many?
        \item Which two queries collide if $s = 0110$? If $s = 1100$?
        \item How many queries for $n = 20$, against the obvious deterministic algorithm and the
              quantum one?
    \end{enumerate}
\end{pproblem}
\tcblower
\begin{wsolution}
    The algorithm splits $x$ into its first $m = \ceil{\sfrac{n}{2}}$ bits $u$ and its last $n - m$
    bits $v$, and queries every string of the form $(u, 0 \dots 0)$ and every string of the form
    $(0 \dots 0, v)$.

    (a) Here $m = 2$. The strings $u00$ are $0000, 0100, 1000, 1100$. The strings $00v$ are
    $0000, 0001, 0010, 0011$. The string $0000$ is in both, so there are $4 + 4 - 1 = 7$ queries.

    (b) $s = 0110$ has $u = 01$ and $v = 10$. The queries $0100$ and $0010$ have XOR $0110 = s$, so
    $f(0100) = f(0010)$. $s = 1100$ has $u = 11$ and $v = 00$. The queries $1100$ and $0000$ have
    XOR $1100 = s$, so $f(1100) = f(0000)$.

    In general every $s = (u, v) \neq 0000$ is the XOR of the two queried strings $(u, 00)$ and
    $(00, v)$, and these are different because $s \neq 0000$. So a collision always appears, and by
    the promise any collision $x, x'$ has $x \oplus x' = s$. If no two answers agree, $f$ is
    one-to-one and $s = 0000$.

    (c) The count is $2^m + 2^{n-m} - 1$. For $n = 20$ that is $1024 + 1024 - 1 = 2047$. Querying
    inputs until two answers agree needs $2^{19} + 1 = 524{,}289$ queries in the worst case. Simon's
    algorithm makes $n + 1 = 21$.

    So the deterministic count is about $2\sqrt{N}$, the same order as the randomized lower bound of
    \cref{sec:lower}, and the quantum count is about $\log N$.
\end{wsolution}
\end{exercisebox}

\begin{exercisebox}
\begin{pproblem}[Grover when $a$ is unknown]\label{pp:unknown-a}
    Let $N = 64$, and let the number of solutions $a$ be unknown, with $1 \leq a \leq 32$.
    \begin{enumerate}[label=(\alph*)]
        \item Why not just run $6$ iterations?
        \item State the algorithm of the Oct 8 exercise and its guarantee.
        \item Compute its success probability for $a = 1$ and for $a = 4$.
        \item How do we make the failure probability small?
    \end{enumerate}
\end{pproblem}
\tcblower
\begin{wsolution}
    (a) Six iterations is right for $a = 1$ (success $0.997$), but for $a = 4$ it rotates past the
    target and succeeds with probability $0.02$ (\cref{pp:wrong-a}). The right $k$ depends on $a$.

    (b) Let $K = \floor{2\sqrt{N}} = 16$. Choose $k$ uniformly from $\set{1, \dots, 16}$, make $k$
    Grover iterations on the uniform state and measure. The exercise proves
    \[\prob{\text{success}}
        = \frac{1}{K} \sum_{k=1}^{K} \sin^2\left((2k+1)\theta\right)
        = \frac{1}{2} - \frac{\sin\left(4(K+1)\theta\right) - \sin 4\theta}{4K \sin 2\theta}
        \geq \frac{1}{2} - \frac{1}{4K \sin 2\theta}
        \geq \frac{3}{8}.\]

    (c) For $a = 1$: $\sin\theta = \sfrac{1}{8}$ and $\theta \approx 0.1253$. Then
    $\sin 2\theta \approx 0.2480$, $\sin 4\theta \approx 0.4806$ and
    $\sin(68\theta) = \sin(8.522) \approx 0.7849$. The probability is
    \[\frac{1}{2} - \frac{0.7849 - 0.4806}{64 \cdot 0.2480} \approx \frac{1}{2} - 0.019 = 0.481.\]
    For $a = 4$: $\sin\theta = \frac{1}{4}$ and $\theta \approx 0.2527$. Then
    $\sin 2\theta \approx 0.4841$, $\sin 4\theta \approx 0.8472$ and
    $\sin(68\theta) = \sin(17.182) \approx -0.9953$. The probability is
    \[\frac{1}{2} - \frac{-0.9953 - 0.8472}{64 \cdot 0.4841} \approx \frac{1}{2} + 0.059 = 0.559.\]
    Both are above $\sfrac{3}{8} = 0.375$. The reason it works: averaged over many values of $k$,
    $\sin^2$ of the angle is about $\frac{1}{2}$, whatever $\theta$ is. The last inequality of (b)
    holds because
    $K \sin 2\theta \geq \sqrt{2N} \cdot \sqrt{\sfrac{2a}{N}} = 2\sqrt{a} \geq 2$, so the subtracted
    term is at most $\sfrac{1}{8}$. It uses $a \leq \sfrac{N}{2}$.

    (d) Check the measured $x$ with one query to $f$. If it is not a solution, repeat with a new
    random $k$. After $r$ rounds the failure probability is at most
    $\left(\frac{5}{8}\right)^r$, which is below $1\%$ for $r = 10$.

    Each round costs at most $16 = 2\sqrt{N}$ queries, so the algorithm is still
    $\bigO{\sqrt{N}}$ without knowing $a$.
\end{wsolution}
\end{exercisebox}


\end{document}
